% CICIDS2018 ML 분류 성능 비교표
% 사용법: \usepackage{booktabs} 필요

\begin{table}[htbp]
\centering
\caption{CICIDS2018 ML 분류 성능 비교 (Binary F1, \%)}
\label{tab:ml_comparison}
\begin{tabular}{lcccccc}
\toprule
방법 & 피처 수 & XGBoost & RF & CNN & LSTM & LogReg \\
\midrule
CASS (제안) & 16 & 93.68 & 93.12 & 90.51 & 92.17 & \textbf{74.29} \\
umar2024~\cite{umar2024} & 12 & \textbf{93.75} & \textbf{93.43} & \textbf{91.58} & \textbf{93.37} & 72.51 \\
\bottomrule
\end{tabular}
\begin{tablenotes}
  \small
  \item CASS는 BM 최대화를 목적함수로 선택된 피처 조합이며,
        umar2024는 분류 성능 기반으로 선택된 12개 피처 조합이다.
        두 방법의 피처 수가 상이함을 고려하여 해석해야 한다.
\end{tablenotes}
\end{table}
